\clearpage\thispagestyle{otFront}\phantomsection\label{avant-propos}\pdfbookmark[0]{Apprendre et comprendre}{bm-avant-propos}\otPageTitre{Apprendre et comprendre}{Un livre pour toute la Quatrième · Nouveau programme}\fontsize{11}{15}\selectfont

Ce livre accompagne une année de mathématiques en Quatrième, selon le
programme du cycle 4 publié en mars 2026 : calculer avec les relatifs et les
fractions, découvrir les puissances, raisonner en géométrie et résoudre des
problèmes qui relient plusieurs idées. Il peut être suivi dans l’ordre ou
utilisé pour reprendre une difficulté précise.\par\medskip
\begin{savoirfaire}[Une progression régulière]Chaque chapitre contient des objectifs, des prérequis, deux pages de cours et quinze exercices : cinq pour découvrir, cinq pour s’entraîner et cinq problèmes. Les corrections expliquent le choix de méthode, les étapes et l’interprétation du résultat.\end{savoirfaire}
\begin{savoirfaire}[Raisonner et démontrer]Le programme met l’accent sur la preuve : définition, propriété, réciproque et contraposée sont distinguées ; les théorèmes de la droite des milieux, de Pythagore et du cercle circonscrit sont démontrés ou expliqués. Le dessin aide à chercher ; les données et les propriétés fondent la preuve.\end{savoirfaire}
\begin{savoirfaire}[Chercher davantage]Dix-huit défis facultatifs proposent une recherche plus longue, avec deux indications séparées et une solution. Deux bilans annuels mélangent les notions. Les fiches outils et le repérage final orientent les révisions.\end{savoirfaire}
Les six compétences sont travaillées : chercher, modéliser, représenter,
raisonner, calculer et communiquer. Le progrès se mesure aussi à la capacité
d’expliquer une démarche et de corriger une erreur.

\clearpage
